\documentclass[12pt,a4paper]{article}

\usepackage{amsmath, amsthm, amssymb, amsfonts}
\usepackage{mathrsfs}
\usepackage{hyperref}
\usepackage{geometry}
\usepackage{enumitem}
\usepackage{booktabs}

\geometry{margin=2.5cm}

\newtheorem{theorem}{Theorem}[section]
\newtheorem{proposition}[theorem]{Proposition}
\newtheorem{lemma}[theorem]{Lemma}
\newtheorem{corollary}[theorem]{Corollary}
\newtheorem{definition}[theorem]{Definition}
\newtheorem{remark}[theorem]{Remark}
\newtheorem{openproblem}[theorem]{Open Problem}
\newtheorem{conjecture}[theorem]{Conjecture}

\newcommand{\lam}[1]{\lambda_{#1}^{(c)}}
\newcommand{\pswf}[1]{\psi_{#1}^{(c)}}
\newcommand{\Tc}{T_c}
\newcommand{\VN}{V_N}
\newcommand{\GN}{G_N}
\newcommand{\inner}[2]{\langle #1,\, #2 \rangle}
\newcommand{\norm}[1]{\left\|#1\right\|}

\title{%
  \textbf{Paper~XII: Direct Coercivity}\\[0.4em]
  \large $\langle (I - T_c)f,\, f \rangle \geq C c^{-1/2} \|f\|^2$
  on $V_N$ without Spectral Factorisation
}
\author{\textit{prolate-gram-coercivity programme}}
\date{May 2026}

\begin{document}
\maketitle

\begin{abstract}
Papers~IX--XI established that Assumption B-strong is
equivalent to a spectral gap $\varepsilon_N \geq Cc^{-1/2}$,
and that neither local (WKB/ORX) nor global (Fredholm
determinant) methods can reach this bound:
the Fredholm determinant is bulk-dominated and blind
to the Galerkin scale, while local methods are
constrained by tunneling to exponential bounds.

\medskip\noindent
This paper pursues a third route:
\emph{direct coercivity} of the quadratic form
\[
  \boxed{
  \inner{(I - \Tc)f}{f} \geq C c^{-1/2} \norm{f}^2
  \quad \forall\, f \in \VN,}
  \tag{Coer}
\]
without passing through eigenvalues $\varepsilon_k$
as intermediate objects.
We formulate three proof routes
(algebraic Gram structure, kernel localisation,
phase-space filling),
identify the precise obstacle in each,
and formulate the minimal open statement
that would close the programme.
\end{abstract}

\tableofcontents
\bigskip\hrule\bigskip

%% ============================================================
\section{Setup and Standing Context}
\label{sec:setup}
%% ============================================================

\subsection{What has been established}

The programme has reduced everything to a single
scalar inequality (Paper~IX, Proposition~2.3):
\[
  \varepsilon_N^{(c)} = 1 - \lam{N} \geq C c^{-1/2},
  \qquad N = \lfloor A c^{1/3} \rfloor.
\]
Papers~X--XI showed that this cannot be reached via:
\begin{itemize}[nosep]
  \item WKB/ORX outer-tail: gives only
        $\varepsilon_N \gtrsim e^{-Cc^{2/3}}$
  \item Fredholm determinant $D_c(1)$: bulk-dominated,
        Galerkin sector below resolution threshold
\end{itemize}

\subsection{The new object}

Instead of asking ``how large is $\varepsilon_N$?''
we ask directly:
\[
  \text{Is } (I - \Tc)|_{\VN}
  \text{ coercive with constant } \geq C c^{-1/2}?
\]
These are equivalent by Paper~IX (since $\VN$ is
$\Tc$-invariant and $\Tc|_{\VN}$ is diagonal in
the PSWF basis).
The advantage of the direct formulation is that
it does not require knowing $\varepsilon_N$
individually --- it requires only a
\emph{lower bound on the quadratic form}
over all $f \in \VN$ simultaneously.

\begin{remark}[Why this is different]
In the eigenvalue formulation, one needs to control
the \emph{smallest} eigenvalue of $(I-\Tc)|_{\VN}$,
which is $\varepsilon_N$.
In the quadratic form formulation, one needs a
\emph{global} lower bound that could in principle
come from the collective structure of $\VN$
rather than from the individual mode $\pswf{N}$.
This is the sense in which (Coer) is a genuinely
different question.
\end{remark}

%% ============================================================
\section{Route A: Algebraic Gram Structure}
\label{sec:gram}
%% ============================================================

\subsection{The Gram matrix}

Let $\{e_k\}_{k=0}^{N}$ be any orthonormal basis
of $\VN$ (e.g.\ the PSWFs $\pswf{k}$).
Define the \emph{prolate Gram matrix}:
\[
  [\GN]_{kl} := \inner{(I-\Tc)e_k}{e_l}
  = \delta_{kl} - \inner{\Tc e_k}{e_l}.
\]
In the PSWF basis, $\GN = \mathrm{diag}(\varepsilon_0,
\ldots, \varepsilon_N)$, so (Coer) is equivalent to
$\lambda_{\min}(\GN) \geq Cc^{-1/2}$.

\subsection{The obstacle in the algebraic route}

In an arbitrary basis, $\GN$ is a full matrix
and its smallest eigenvalue need not be
directly related to $\varepsilon_N$.
However, in any basis, the Courant--Fischer theorem gives:
\[
  \lambda_{\min}(\GN)
  = \min_{\norm{v}=1} v^\top \GN v
  = \min_{f \in \VN,\, \norm{f}=1}
    \inner{(I-\Tc)f}{f}.
\]
Since $\VN$ is $\Tc$-invariant,
this minimum is always achieved by $f = \pswf{N}$,
giving $\lambda_{\min}(\GN) = \varepsilon_N$.
The algebraic route therefore \emph{reduces back to}
the eigenvalue problem.

\begin{proposition}[Algebraic route collapses]
\label{prop:gram_collapse}
For any choice of basis of $\VN$:
\[
  \lambda_{\min}(\GN) = \varepsilon_N^{(c)}.
\]
The algebraic Gram route does not bypass the
eigenvalue problem; it is equivalent to it.
\end{proposition}

\begin{remark}[What the algebraic route does give]
Proposition~\ref{prop:gram_collapse} is not a
dead end: it precisely locates the difficulty.
The coercivity constant is $\varepsilon_N$,
achieved by the single mode $\pswf{N}$.
To prove (Coer), one must therefore either:
\begin{enumerate}[nosep]
  \item show $\varepsilon_N \geq Cc^{-1/2}$ directly, or
  \item show that the minimum of
        $\inner{(I-\Tc)f}{f}$ over $\VN$
        is \emph{not} achieved by a pure mode
        (which would require $\VN$ to not be
        $\Tc$-invariant, contradicting the PSWF structure).
\end{enumerate}
Option (2) is impossible for the standard $\VN
= \mathrm{span}\{\pswf{0},\ldots,\pswf{N}\}$.
It becomes available only if $\VN$ is replaced by
a \emph{different} $N$-dimensional subspace of $L^2[-1,1]$
that is not $\Tc$-invariant.
\end{remark}

%% ============================================================
\section{Route B: Kernel Localisation}
\label{sec:kernel}
%% ============================================================

\subsection{The quadratic form as a double integral}

For $f \in \VN$ with $\norm{f} = 1$:
\[
  \inner{(I-\Tc)f}{f}
  = \int_{-1}^{1} |f(x)|^2\,dx
  - \int_{-1}^{1}\int_{-1}^{1}
    \frac{\sin(c(x-y))}{\pi(x-y)}\,
    f(x)\overline{f(y)}\,dx\,dy.
\]
Since $\Tc f = \sum_{k=0}^N \lam{k}
\inner{f}{\pswf{k}}\pswf{k}$:
\[
  \inner{\Tc f}{f}
  = \sum_{k=0}^N \lam{k} |\hat f_k|^2,
  \quad \hat f_k := \inner{f}{\pswf{k}},
\]
so
\[
  \inner{(I-\Tc)f}{f}
  = \sum_{k=0}^N (1-\lam{k})|\hat f_k|^2
  = \sum_{k=0}^N \varepsilon_k |\hat f_k|^2.
\]
This is a weighted $\ell^2$-norm of the PSWF
coefficients of $f$, with weights $\varepsilon_k$.

\subsection{The obstacle in the kernel route}

The minimum of $\sum_k \varepsilon_k |\hat f_k|^2$
subject to $\sum_k |\hat f_k|^2 = 1$ is again
$\min_k \varepsilon_k = \varepsilon_N$.
The kernel route therefore also collapses to
the eigenvalue problem.

\begin{proposition}[Kernel route collapses]
\label{prop:kernel_collapse}
\[
  \min_{f \in \VN,\, \norm{f}=1}
  \inner{(I-\Tc)f}{f}
  = \varepsilon_N^{(c)}.
\]
\end{proposition}

\begin{remark}[What the kernel route does give --- a positive direction]
\label{rem:kernel_positive}
Although the minimum is $\varepsilon_N$,
the kernel localisation reveals a useful structure:
if $f$ is \emph{not} aligned with $\pswf{N}$,
i.e.\ $|\hat f_N|^2 \leq 1 - \delta$ for some $\delta > 0$,
then:
\[
  \inner{(I-\Tc)f}{f}
  \geq \varepsilon_{N-1}(1-\delta) + \varepsilon_N \delta
  \geq \varepsilon_{N-1}(1-\delta).
\]
For $k < N$, $\varepsilon_k \gg \varepsilon_N$
(since eigenvalues are monotone), so
the form is much larger than $\varepsilon_N$
unless $f$ is nearly $\pswf{N}$.

\medskip\noindent
This suggests an \emph{energy separation argument}:
if one can show that the relevant functions
$f$ in the Galerkin approximation have
$|\hat f_N|^2 \leq 1 - c^{-\alpha}$
for some $\alpha < 1/2$, then (Coer)
follows from $\varepsilon_{N-1}$ alone,
without needing to control $\varepsilon_N$.
\end{remark}

%% ============================================================
\section{Route C: Phase-Space Filling}
\label{sec:phasespace}
%% ============================================================

\subsection{The underfilling argument}

For $k \leq N = O(c^{1/3})$, the PSWF $\pswf{k}$
has approximately $k$ zeros on $[-1,1]$ and
Fourier support essentially in $[-c,c]$.
Its phase-space ``footprint'' is a rectangle of
area $\sim 4c \cdot 2 = 8c$.
But $k \leq c^{1/3}$ modes cannot fill this
rectangle efficiently: each mode occupies
a phase-space area $\sim 2\pi$ (by the
uncertainty principle), so the $N$ lowest modes
fill an area $\sim 2\pi N \sim c^{1/3}$,
leaving a fraction $1 - c^{1/3}/(8c) \approx 1$
of the phase-space rectangle unfilled.

\begin{proposition}[Phase-space underfilling]
\label{prop:phasespace}
For $f \in \VN$ with $\norm{f}=1$:
the Wigner distribution $W_f(x,\xi)$ has
$L^2$-mass concentrated in a region of
phase-space area $O(c^{1/3})$,
while $\Tc$ acts as a projection onto a
phase-space region of area $\sim 4c \cdot 2$.
The \emph{uncovered fraction} is
$1 - O(c^{-2/3})$, i.e.\ almost the entire
phase-space rectangle.
\end{proposition}

\subsection{Why underfilling gives coercivity --- heuristically}

The energy ``leaked'' by $f$ outside $[-1,1]$
is precisely $\inner{(I-\Tc)f}{f}$.
Phase-space underfilling says that $f$'s
energy is concentrated in a small region
of phase-space.
If this region is generically not aligned with the
``exit channels'' of $\Tc$ (the directions of
maximal concentration loss), then the leaked
energy is at least the fraction
of phase-space that is ``underfilled'' ---
which is $\Omega(1)$.

\begin{remark}[Why this heuristic fails to give $c^{-1/2}$]
The phase-space underfilling argument gives a
constant lower bound (independent of $c$),
not a $c^{-1/2}$ bound.
The $c^{-1/2}$ scale must come from a
\emph{quantitative} version of the underfilling,
specifically from the $O(c^{1/3})$ transition
window in the Airy regime:
the ``exit channels'' of $\Tc$ are localised
in a phase-space strip of width $O(c^{-1/3})$
(in the spectral parameter), and the overlap
of $f \in \VN$ with this strip is
$O(c^{1/3}) \cdot c^{-1/3} = O(1)$.
Turning this into a $c^{-1/2}$ bound
requires a more precise estimate of the
overlap, which is the content of the
open problem below.
\end{remark}

%% ============================================================
\section{The Minimal Open Statement}
\label{sec:open}
%% ============================================================

\begin{openproblem}[Minimal statement for (Coer)]
\label{op:minimal}
Show that for all $f \in \VN$
with $\norm{f} = 1$ and
$N = \lfloor Ac^{1/3} \rfloor$:
\[
  |\inner{f}{\pswf{N}}|^2
  \leq 1 - \delta(c)
\]
for some $\delta(c) \geq C c^{-1/2}$,
i.e.\ no unit vector in $\VN$
is more than $1 - Cc^{-1/2}$
aligned with the worst mode $\pswf{N}$.

\medskip\noindent
By Remark~\ref{rem:kernel_positive},
this would imply (Coer) with the constant
coming from $\varepsilon_{N-1}$ rather than $\varepsilon_N$.
\end{openproblem}

\begin{remark}[Why Open Problem~\ref{op:minimal} is likely false for $\VN$]
In the PSWF basis, $\pswf{N} \in \VN$ by definition,
so $|\inner{\pswf{N}}{\pswf{N}}|^2 = 1$.
The minimum alignment is achieved by $f = \pswf{N}$
itself, giving $\delta = 0$.
Open Problem~\ref{op:minimal} is therefore
\emph{vacuously false} for the standard $\VN$.

\medskip\noindent
This reinforces the conclusion of
Proposition~\ref{prop:gram_collapse}:
for the PSWF-defined $\VN$, the direct
coercivity approach is equivalent to the
eigenvalue approach.
(Coer) holds if and only if $\varepsilon_N \geq Cc^{-1/2}$.
\end{remark}

\begin{openproblem}[Alternative subspace]
\label{op:altspace}
Define $\widetilde{V}_N$ to be an $N$-dimensional
subspace of $L^2[-1,1]$ that is \emph{not}
$\Tc$-invariant but approximates $\VN$ in the
sense that the Galerkin approximation in
$\widetilde{V}_N$ has the same order of accuracy
as in $\VN$.
Does there exist such a $\widetilde{V}_N$ for which:
\[
  \min_{f \in \widetilde{V}_N,\, \norm{f}=1}
  \inner{(I-\Tc)f}{f}
  \geq C c^{-1/2}?
\]
If so, (Coer) holds for $\widetilde{V}_N$,
and the Paper~VI error estimate may be recoverable
with $\widetilde{V}_N$ replacing $\VN$.
\end{openproblem}

%% ============================================================
\section{Option B: Phase-Space Coercivity as Escape Hatch}
\label{sec:optionB}
%% ============================================================

The preceding analysis shows that (Coer)
for the standard $\VN$ is equivalent to
$\varepsilon_N \geq Cc^{-1/2}$,
and that all three direct routes collapse
to this eigenvalue bound.

\medskip\noindent
The structural reason is that $\VN$ is
$\Tc$-invariant by definition.
The only route that could bypass this is to
replace $\VN$ by a subspace that is
\emph{not} $\Tc$-invariant
(Open Problem~\ref{op:altspace}).

\medskip\noindent
This is precisely the setting of
\textbf{microlocal / phase-space coercivity}
(Option B of the programme):
instead of the PSWF eigenspace $\VN$,
use a \emph{phase-space cutoff space}
$\widetilde{V}_N = \Pi_{P} L^2$,
where $\Pi_P$ is a pseudodifferential
cutoff to a phase-space region $P \subset
\mathbb{R} \times \mathbb{R}$ of area $\sim N$.

\begin{remark}[Why Option B is the escape hatch]
In the phase-space cutoff framework:
\begin{itemize}[nosep]
  \item $\widetilde{V}_N$ is \emph{not} $\Tc$-invariant
        (it is a spatial cutoff, not a spectral one).
  \item The coercivity of $(I-\Tc)$ on $\widetilde{V}_N$
        is controlled by the \emph{overlap} of the
        phase-space region $P$ with the
        ``concentration zone'' of $\Tc$
        (i.e.\ $[-1,1] \times [-c,c]$).
  \item For $P$ chosen as a phase-space rectangle
        of area $\sim c^{1/3}$ aligned with the
        low-frequency modes, the overlap is
        $O(c^{1/3}/c) = O(c^{-2/3})$,
        giving $(I-\Tc)$ coercive with constant
        $\sim 1 - O(c^{-2/3}) \approx 1$
        --- much better than $\varepsilon_N$.
\end{itemize}
The cost is that $\widetilde{V}_N \neq \VN$,
so the Galerkin approximation in $\widetilde{V}_N$
requires separate accuracy analysis.
This is the content of Paper~XIII (if needed).
\end{remark}

\begin{openproblem}[Phase-space subspace approximation]
\label{op:phasespace_approx}
Let $\widetilde{V}_N = \Pi_P L^2[-1,1]$ for
$P = [-1,1] \times [-c^{1/3}, c^{1/3}]$.
Show that:
\begin{enumerate}[nosep]
  \item $(I-\Tc)$ is coercive on $\widetilde{V}_N$
        with constant $\geq Cc^{-1/2}$.
  \item The best approximation error in
        $\widetilde{V}_N$ is of the same order
        as in $\VN$ for the target function class.
\end{enumerate}
Together, (1) and (2) would replace B-strong
entirely with a phase-space coercivity statement.
\end{openproblem}

%% ============================================================
\section{Summary and Programme State}
\label{sec:summary}
%% ============================================================

\begin{center}
\renewcommand{\arraystretch}{1.5}
\begin{tabular}{p{2.2cm} p{4.5cm} p{4.5cm} p{1.5cm}}
\toprule
\textbf{Route} & \textbf{Idea} & \textbf{Outcome} & \textbf{Status} \\
\midrule
A (Gram) &
  Direct $\lambda_{\min}(G_N) \geq Cc^{-1/2}$ &
  Collapses to $\varepsilon_N$ (Prop.~\ref{prop:gram_collapse}) &
  Closed \\
B (Kernel) &
  Weighted $\ell^2$ form over modes &
  Collapses to $\varepsilon_N$ (Prop.~\ref{prop:kernel_collapse});
  energy separation possible if $f \not\approx \psi_N$ &
  Partial \\
C (Phase-space) &
  Underfilling of phase-space rectangle &
  Gives $O(1)$ bound, not $c^{-1/2}$; needs quantitative overlap &
  Open \\
\midrule
\textbf{Option B} &
  Replace $V_N$ by phase-space cutoff $\widetilde{V}_N$ &
  Coercivity likely $\geq Cc^{-1/2}$; accuracy TBD &
  \textbf{Open} \\
\bottomrule
\end{tabular}
\end{center}

\medskip\noindent
\textbf{Programme state after Paper~XII.}
The direct coercivity approach establishes that
(Coer) for the standard $\VN$ is rigorously
equivalent to $\varepsilon_N \geq Cc^{-1/2}$
--- confirming that no shortcut through
Gram structure or kernel localisation exists
for the PSWF eigenspace.
The only genuine bypass is the phase-space
subspace replacement of Open Problem~\ref{op:phasespace_approx},
which is the starting point of Paper~XIII.

%% ============================================================
\begin{thebibliography}{9}

\bibitem{paper9}
  \textit{Paper~IX: B-strong as a Spectral Gap},
  \texttt{paper9\_superconcentration.tex},
  prolate-gram-coercivity programme, 2026.

\bibitem{paper11}
  \textit{Paper~XI: Fredholm Determinant Microlocal Inversion},
  \texttt{paper11\_fredholm\_microlocal.tex},
  prolate-gram-coercivity programme, 2026.

\bibitem{LP1961II}
  H.~J.~Landau and H.~O.~Pollak,
  ``Prolate Spheroidal Wave Functions,
  Fourier Analysis and Uncertainty~II,''
  \textit{Bell Syst.\ Tech.\ J.}, 40 (1961), 65--84.

\bibitem{ORX2013}
  A.~Osipov, V.~Rokhlin, H.~Xiao,
  \textit{Prolate Spheroidal Wave Functions of Order Zero},
  Springer, 2013.

\bibitem{Folland1989}
  G.~B.~Folland,
  \textit{Harmonic Analysis in Phase Space},
  Princeton University Press, 1989.

\end{thebibliography}

\end{document}
